\documentclass[10pt,a4paper]{amsart}
\usepackage[margin=19mm]{geometry}
\usepackage{amsmath,amssymb,mathtools}
\usepackage[scr=rsfs,cal=euler]{mathalfa}
\usepackage{xfrac}
\usepackage[unicode,hidelinks,bookmarks=false]{hyperref}
\newcommand{\ie}{i.e.\ }
\newcommand{\cf}{cf.\ }
\newcommand{\resp}{respectively\ }
\newcommand{\Subsection}{\subsection}
\providecommand{\nobreakdash}{\nobreak\hbox{-}}
\setlength{\emergencystretch}{2em}
\multlinegap0pt
\usepackage{amssymb}
\usepackage{enumerate}
\usepackage{algorithm}\usepackage{algpseudocode}
\usepackage{mathtools}
\usepackage{tikz}
\newtheorem{theorem}{Theorem}
\newtheorem{claim}{Claim}
\title{Three-coloring triangle-free graphs without long forbidden paths}
\author{Yidong Zhou, Jorik Jooken, Baoyuan Shan, Jan Goedgebeur, Shenwei Huang}
\date{Source: arXiv:2512.12349v2 (CC BY 4.0)}
\begin{document}
\maketitle

\noindent\emph{Source and licence.} Three-coloring triangle-free graphs without long forbidden paths. Version arXiv:2512.12349v2 (CC BY 4.0), distributed by arXiv under the Creative Commons Attribution 4.0 International licence, http://creativecommons.org/licenses/by/4.0/, as recorded in the arXiv licence metadata for this exact version and shown on the abstract page https://arxiv.org/abs/2512.12349v2. Full licence notice: \texttt{licenses/arxiv-2512.12349-CC-BY-4.0.txt}.\par\smallskip

\noindent\textit{Editorial scope.} Counted unit: the article's theorem P11 (source0.tex lines 759--761). The article's complete original proof of it is delivered with it (source0.tex lines 762--838, one \texttt{proof} environment of 77 lines). No mathematics is written, repaired or re-derived: the statement and the proof are the article's own lines, in the article's own notation, and no retained line is substituted or re-flowed.  No further source-local context is needed: every cross-reference of every retained line resolves inside the two delivered spans.  Nothing was omitted from any retained span, so the package carries no omission record.  No retained line contains a citation, so the wrapper carries no bibliography and no bibliography entry.  No retained line contains a figure environment, an image-inclusion command or drawing code, so no image is delivered and the package declares no asset dependency.  Editorial formatting only, in addition: an AMS \texttt{amsart} preamble that restates the article's own declarations and macros used by the retained lines, namely the environments \texttt{theorem}, \texttt{claim} on separate counters, together with the standard packages those lines need.  The retained environments print in this document as Theorem 1, Claim 1 in order of appearance, whereas the article prints the same items with the numbering of its own full document.  The complete native source of the article as distributed in the arXiv e-print for this exact version is bundled unchanged as \texttt{sources/190937/source0.tex}.
\begin{theorem}\label{P11}
    For each $k\geq 3$, $G_k$ is a 4-vertex-critical $\{4K_2,C_3\}$-free graph. 
\end{theorem}

\begin{proof}
    We start our proof by showing that $G_k$ is $\{4K_2,C_3\}$-free. 
    \begin{claim}\label{claim:P_11 C_3-free}
        For every $k\geq 3$, $G_k$ is $\{4K_2,C_3\}$-free. 
    \end{claim}
    \begin{proof}
        $G_k$ is clearly $C_3$-free because of the definition of $G_k$. Therefore, we only need to prove that $G_k$ is $4K_2$-free. 
        Suppose for the sake of obtaining a contradiction that $G_k$ contains an induced $4K_2$, say $H$.  
        We suppose first that $v_iv_{i+1}$ is an edge of $H$ for some $i$.  
        Note that $M(v_i,v_{i+1})=\{v_{i+3}\}\cup \{v_{i+4},v_{i+7},v_{i+10},\ldots, v_{i+3k+4}, v_{i+3k+7}\}\cup \{v_{i-2}\}$. 
        Since $H$ is a $4K_2$, $G[M(v_i,v_{i+1})]$ contains a $3K_2$. 
        However, note that $M(v_i,v_{i+1}) \setminus \{v_{i+3},v_{i-2}\} =\{v_{i+4},v_{i+7},v_{i+10},\ldots, v_{i+3k+4}, v_{i+3k+7}\}$ is a stable set and therefore $G[M(v_i,v_{i+1})]$ is $3K_2$-free, a contradiction. So we may assume that for any $i\in [3k+10]$, $v_iv_{i+1}$ is not an edge of $H$. 
        
        Suppose that $v_iv_{i+5+3j}$ is an edge of $H$, where $0\leq j\leq k$. 
        Note that  
	\begin{equation}
		\begin{aligned}
			M(v_i,v_{i+5+3j})=\{v_{i+2}\}\cup A\cup \{v_{i+3+3j}\} \cup \{v_{i+7+3j}\}\cup B\cup \{v_{i-2}\},
		\end{aligned}
	\end{equation}
	where $$A=\{v_{i+4},v_{i+7},v_{i+10},\ldots, v_{i+1+3j}\}, B=\{v_{i+9+3j},v_{i+12+3j},\ldots, v_{i+9+3(k-1)}\}.$$ 
	Note that $A=\emptyset$ if $j=0$ and $B=\emptyset$ if $j=k$. 
    Since $A$ and $B$ are both stable sets and $A$ is complete to $B$, $|V(H)\cap (A\cup B)|\leq 3$. 
	Thus, $|V(H)\cap \{v_{i+2},v_{i+3+3j},v_{i+7+3j},v_{i-2}\}|\geq 3$. 
	So by symmetry, we may assume that $v_{i+2},v_{i+7+3j}\in V(H)$. Note that $v_{i+2}$ is complete to $A\setminus \{v_{i+4}\}$ and $v_{i+7+3j}$ is complete to $B\setminus \{v_{i+9+3j}\}$. 
    It follows that $V(H)\cap (A\cup B)\subseteq \{v_{i+4},v_{i+9+3j}\}$. 
    On the other hand, since $|V(H)\cap M(v_i,v_{i+5+3j})|=6$, $|V(H)\cap (A\cup B)|\geq 2$. 
    Thus, $V(H)\cap (A\cup B)=\{v_{i+4},v_{i+9+3j}\}$. 
    So $V(H)\cap M(v_i,v_{i+5+3j})=\{v_{i+2}, v_{i+4}, v_{i+3+3j}, v_{i+7+3j}, v_{i+9+3j}, v_{i-2}\}$. 
    Note that $v_{i-2}v_{i+3+3j},v_{i+4}v_{i+9+3j}\in E(G_k)$. 
    Since $k\geq 3$, either $j\neq 1$ or $j\neq k-1$. 
    If $j\neq 1$, then $v_{i+3+3j}v_{i+4}\in E(G_k)$. 
    If $j\neq k-1$, then $v_{i-2}v_{i+9+3j}\in E(G_k)$. 
    Each case contradicts that $H$ is an induced $4K_2$. 
    This completes the proof of Claim \ref{claim:P_11 C_3-free}.
    \end{proof}

    Next, we show that every graph obtained by deleting one vertex from $G_k$ admits a proper 3-coloring. 
    \begin{claim}\label{claim:critical of G_k}
        For every $v\in V(G_k)$, $\chi(G_k\setminus \{v\})\leq 3$. 
    \end{claim}
    \begin{proof}
        Let $c: V(G_k\setminus \{v_1\}) \to \{0,1,2\}$ be a function such that $c(v_i)\equiv i\pmod{3}$. 
        We then prove that $c$ is a proper 3-coloring of $G_k\setminus \{v_1\}$. 
        Suppose to the contrary that there are two adjacent vertices $v_i,v_j\in V(G_k\setminus \{v_1\})$ with $i<j$ such that $c(v_i)=c(v_j)$.
        So $j-i \equiv 0\pmod{3}$. 
        By the definition of $G_k$, since $v_i$ and $v_j$ are adjacent, $j=i+1$ or $j=i+5+3\ell$ for some integer $\ell \geq 0$. 
        Hence, $j-i \equiv 1\pmod{3}$ or $j-i \equiv 2\pmod{3}$, which is a contradiction. 
        By the symmetry of $G_k$, for every $v\in V(G)$, $\chi(G\setminus \{v\})\leq 3$. 
        This completes the proof of Claim \ref{claim:critical of G_k}.
    \end{proof}
    By Claim \ref{claim:critical of G_k}, $\chi(G_k)\leq 4$. 
    Suppose now that there is a proper 3-coloring $c$ of $G_k$. 
	Since $\{v_1,v_2,v_3,v_4,v_5\}$ induces a $C_5$, $\chi(G_k)\geq 3$. 
	So there exists an index $i$ such that the colors of $v_i,v_{i+1},v_{i+2}$ are pairwise different. 
	By symmetry, we may assume that $c(v_i)=1,c(v_{i+1})=2$ and $c(v_{i+2})=3$. 
    \begin{claim}\label{claim:3-coloring contradict}
        The colors of $v_{i+1},v_{i+2},v_{i+3}$ are also pairwise different. 
        In particular, $c(v_{i+3})=c(v_i)$.
    \end{claim}
   \begin{proof}
       Suppose to the contrary that $c(v_{i+3})=2$. 
       So $c(v_{i+8})=3$. 
       If $c(v_{i+4})=1$, then $c(v_{i+9})=3$, which is a contradiction. 
       So we may assume that $c(v_{i+4})=3$. 
       Hence, $c(v_{i+9})=1$.  
       If $c(v_{i-1})=3$, then since $c(v_{i+3})=2$, we have $c(v_{i-2})=1$. 
       Then since $v_{i-2}v_{i+9}\in E(G)$ and both of these two vertices are colored 1, we obtain a contradiction.  
       So we may assume that $c(v_{i-1})=2$. But then $c(v_{i+7})=1$, $c(v_{i+6})=3$ and $c(v_{i-2})=1$, which leads again to a contradiction with $c(v_{i+9})=1$. 
       This completes the proof of Claim \ref{claim:3-coloring contradict}.
   \end{proof}
		
	By Claim \ref{claim:3-coloring contradict}, if $j\equiv i\pmod{3}$, then $c(v_j)=c(v_i)$. 
    Thus, $c(v_{i+3k+9})=c(v_{i})=1$, which is a contradiction. 
    It follows that for every $k\geq 3$, $G_k$ is 4-vertex-critical. 
    This completes the proof of Theorem \ref{P11}.
\end{proof}

\end{document}
